\section{Related Work}
\label{sec:relatedworks}

\subsection{Visual Sim-to-Real and Perceptive WBC}

% Dextrah

% VBC

% Early RGB work in manipulation https://arxiv.org/pdf/1703.06907

% Generative method https://arxiv.org/pdf/2411.00083#page=7.39

Visual sim-to-real has enabled robust visuomotor control across locomotion~\citep{miki2022learning,zhuang2024humanoid,cheng2023parkour,wang2025beamdojo,long2025learning,ren2025vb,xue2025leverb} and manipulation~\citep{sadeghi2016cad2rl,tobin2017domain,peng2018sim,andrychowicz2020learning,akkaya2019solving,handa2023dextreme,yuan2024learning,jiang2024transic,ze2023visual,hansen2021generalization,huang2022spectrum,singh2024dextrah,deng2025graspvla,liu2024visual,zhu2018reinforcement}. For manipulation, domain randomization has long been used to bridge the reality gap for RGB-trained policies~\citep{sadeghi2016cad2rl,tobin2017domain,peng2018sim}. Teacher–student pipelines further improve generalization by distilling privileged policies into visual ones trained under randomized rendering~\citep{singh2024dextrah,deng2025graspvla}; for example, Dextrah-RGB attains zero-shot dexterous grasping from stereo images~\citep{singh2024dextrah}. For whole-body loco-manipulation, VBC integrates legged locomotion and arm control via hierarchical distillation~\citep{liu2024visual}, and generative pipelines~\citep{yu2024learning} scale visual diversity beyond handcrafted assets. LeVERB~\citep{xue2025leverb} takes advantage of photorealistic demos rendered in IsaacSim for coarse whole-body vision-language tasks. Despite this progress, most prior efforts target isolated arms or decouple locomotion from manipulation, and few demonstrate a vision-only, end-to-end policy that solves tasks demanding whole-body capabilities. We address this gap by learning from RGB and proprioception to produce unified loco-manipulation for door opening, without hand-coded primitives or depth/pose priors.


\subsection{Loco-manipulation}

% Some method relies heavily on hard-coded primitives and ground truth perception from LiDAR. Not scalable to real world https://arxiv.org/pdf/2411.03532

% https://icradooropen.github.io/icradooropen/# lacks deliberate manipulation behavior like grasping door handle
% Loco-manipulation requires coordinating base motion and contact-rich manipulation over long horizons while the robot moves through its environment. Sim-to-real efforts span modular designs that decouple legs and arms~\cite{liu2024visual,ben2025homie,cheng2024express} and end-to-end policies for full-body control~\cite{ha2024umi,pan2025roboduet,he2024omnih2o}; some rely on large-scale real data or vision–language models~\cite{qiu2024learning,wang2025odyssey,bjorck2025gr00t}, but robustness remains limited.
% Door opening is a canonical instance of loco-manipulation—approach, stabilize, grasp and actuate the handle, then sweep and traverse. Prior systems typically assume structure or embed hand-engineered primitives, including rule-based pipelines~\cite{oh2017technical,calvert2025behavior}, stagewise controllers without deliberate grasping~\cite{lee2025stageact}, adaptation-heavy designs~\cite{xiong2024adaptive}, and simulator/planner variants that simplify actuation or require privileged sensing~\cite{urakami2019doorgym,traverse-doors-2025,weng2025hdmi}. These choices often target wheeled platforms or reduced-fidelity contact and do not yield an autonomous vision-based policy.
% We treat door opening explicitly as loco-manipulation and train a vision-only, end-to-end controller on a bipedal humanoid that coordinates navigation with manipulation to handle diverse handle types—without brittle primitives or heavy task-specific engineering.

% Most door-opening methods rely on structured environments or hand-engineered primitives. Rule-based pipelines have been used to complete the task~\cite{oh2017technical,calvert2025behavior}. StageACT segments behaviors into stages but does not perform deliberate handle grasping~\cite{lee2025stageact}. DoorBot depends on extensive real-world adaptation~\cite{xiong2024adaptive}. Simulator-driven efforts such as Door-Gym simplify actuation with hooks or assume depth-based priors~\cite{urakami2019doorgym}. Zhang et al.~\cite{traverse-doors-2025} plan with ground-truth localization but offer limited open-world generalization. HDMI~\cite{weng2025hdmi} requires motion-capture hardware, which restricts its use in everyday settings. These approaches typically target wheeled platforms or reduced-fidelity manipulation and do not deliver a general, vision-based policy. In contrast, we train a vision-only, end-to-end controller for a bipedal humanoid that handles diverse handle types, performs coordinated whole-body loco-manipulation, and avoids brittle primitives and heavy task-specific engineering.

Loco-manipulation requires a robot to coordinate whole-body motion, balance, perception, and contact-rich manipulation while navigating through its environment. Sim-to-real efforts span modular designs that decouple legs and arms~\citep{liu2024visual,ben2025homie,cheng2024express} and end-to-end policies for whole-body control~\citep{ha2024umi,pan2025roboduet,he2024omnih2o}. Articulated-object interaction (e.g., doors, drawers, and latches) provides a representative instance of loco-manipulation. Prior systems often embed task-specific structure such as rule-based sequences~\citep{oh2017technical,calvert2025behavior}, stagewise controllers with limited grasp synthesis~\citep{lee2025stageact}, or adaptation-heavy pipelines requiring extensive real-world data~\citep{xiong2024adaptive}. Simulation-driven approaches~\citep{urakami2019doorgym,traverse-doors-2025,weng2025hdmi} frequently assume privileged sensing, simplified actuation, or wheeled platforms, limiting their ability to generalize to unstructured environments and whole-body humanoid control.

\subsection{Reinforcement Learning Fine-Tuning}

% \haoru{Wenli TODO}
RL fine-tuning aims to refine a robot policy through additional interaction with the environment. Generalist agents such as RoboCat~\citep{robocat} interleave rollouts, relabeling, and policy updates to gradually expand manipulation skills and robustness~\citep{robocat}. Self-improving visuomotor systems further show that alternating deployment and learning can close the gap between supervised policies and robust real-world controllers~\citep{medals,sime}. Many of these methods use reinforcement learning to adapt an imitation-learned policy to real-world dynamics~\citep{johannink2019residual,ankile2024resip,xiao2025pld}, often requiring substantial additional real-world interaction. Our fine-tuning phase follows this line of work but targets zero-shot sim-to-real transfer for RGB-based humanoid loco-manipulation. After the DAgger distillation phase, we apply GRPO~\citep{grpo2024} in simulation for policy refinement. Combined with extensive domain randomization, this self-refinement yields a policy that keeps the handle in view, maintains stable whole-body balance during traversal, and recovers from off-nominal camera poses.